\BeleskeID{03-s011}
% element: 03-s011:e01
% element: 03-s011:e02
% element: 03-s011:e03
% element: 03-s011:d01

Neka \(m_k\) predstavlja minterm izabrane adrese. Kada je komponenta izabrana, izlaz sa aktivnom nulom ima vrednost \(\overline{m_k}\): izabrani izlaz je nula, a ostali su jedinice. Nadcrt uz naziv izlaza ili negacioni kružić predstavljaju dva alternativna zapisa te negacije. Zapis sa kružićem uz priključak često se koristi. Njihovo istovremeno unošenje u ovim simbolima daje \(\overline{\overline{m_k}}=m_k\), pa izabrani izlaz opet postaje jedinica.

Za prva dva donja simbola komponenta je izabrana kada je \(CS_1=1\) i \(CS_2=0\). U trećem je CS1 negiran, a CS2 označen i nadcrtom i kružićem. Time uslov postaje \(\overline{CS_1}CS_2\), pa komponentu bira kombinacija \(CS_1=0,CS_2=1\). Jednak broj pinova i isti natpis dekodera ne znače jednaku funkciju kada se razlikuju negacije.
